\documentclass{article}
\usepackage[utf8]{inputenc}
\usepackage[T1]{fontenc}
\usepackage{geometry}
\usepackage{graphicx}
\usepackage{float}
\usepackage{hyperref}
\usepackage[english,indonesian]{babel}

\title{Optimasi Pemrosesan Template LaTeX Menggunakan Pendekatan Berbasis AI}
\author{Academic Support Agent}
\date{\today}

\begin{document}
\maketitle

\begin{abstract}

\end{abstract}

\chapter{PENDAHULUAN}


\chapter{TINJAUAN PUSTAKA}


% Original content (commented out):
% \documentclass[12pt]{article}
\usepackage[spanish]{babel}
\usepackage[utf8]{inputenc}
\usepackage[T1]{fontenc}
\usepackage{graphicx}
\usepackage{cite}
\usepackage{float}
\usepackage{amsfonts}
\usepackage{amsmath}
\usepackage{geometry}
\graphicspath{{IMAGES}}
\newtheorem{tab}{Tabla}
\renewcommand\tablename{Tabla}
\setlength{\columnsep}{7mm}
\renewcommand\tablename{Tabla}
\renewcommand\contentsname{\'Indice}
\setlength{\topmargin}{-0.75in} \setlength{\textheight}{9.25in}
\setlength{\oddsidemargin}{0.0in} \setlength{\evensidemargin}{0.0in}
\setlength{\textwidth}{6.5in}


\begin{document}

\title{PRÁCTICA: GLOBO INFLADO}
\author{Toriz R. Ángel David 201628195}


\twocolumn[
\begin{@twocolumnfalse}
\maketitle
\vspace*{-0.8cm}
\begin{center}\rule{0.9\textwidth}{0.1mm} \end{center}
\begin{abstract}
\normalsize Se da respuesta a la pregunta "¿qué pesa más, un globo inflado o uno desinflado?"  por medio del principio de Arquímedes.
\begin{center}\rule{0.9\textwidth}{0.1mm} \end{center}
\vspace*{0.3cm}
\end{abstract}
\end{@twocolumnfalse}
]

\section{Introducci\'on}
Desde la anitgüedad, gracias al matemático Arquímedes, es conocido el principio de flotación o \emph{principio de Arquímedes}, en honor a su descubridor. El principio enuncia lo siguiente:
\begin{center}
    "Sobre cualquier cuerpo sumergido total o parcialmente en un fluido en reposo actúa una fuerza dirigida hacia arriba, la \emph{fuerza de flotación}, cuya magnitud es igual al peso del fluido desplazado por el cuerpo"\cite{pda:1}
\end{center}
En forma de ecuación:
\begin{equation}
    |F_{f}|=m_{d}\cdot g=V_{d}\rho_{f}\cdot g
\end{equation}
donde $F_{f}$ es la fuerza de flotación, $m_{d}$ es la masa del fluido desplazado, $\rho_{f}$ la densidad del fluido, $V_{d}$ el volumen de fluido desplazado, y $g$ la aceleración debida a la gravedad.
En esta práctica, se utiliza este enunciado para resolver la pregunta "¿qué pesa más, un globo inflado o uno desinflado?". Para este fin, se usa una balanza; la balanza de dos platos usada en el experimento funciona de la siguiente manera al comparar el peso de dos objetos: se coloca un objeto en cada plato, y se espera que lleguen a una posición de equilibrio, es decir, una posición en la que la balanza ya no se mueva. Puesto que lo que mide la balanza es la fuerza normal del plato sobre el objeto, o sea, la fuerza de reacción al peso del objeto (menos la fuerza de flotación, pues ambos objetos están rodeados por aire), el plato más inclinado es aquél para el cual la fuerza normal es mayor:
\begin{equation} \label{eq:2}
    F_{N}=F{g}-F_{f}
\end{equation}
En general, puesto que $\rho_{aire}$ es muy pequeña, la fuerza de flotación debida al aire es despreciable, pero puesto que en esta práctica se usan objetos muy ligeros, debe tomarse en consideración.

\subsection{Arreglo Experimental}
Se utilizaron los siguientes materiales:
\begin{itemize}
\item Balanza de dos platos
\item Globos
\item Bolsas de plástico
\item Botellas de plástico, medianas y grandes
\end{itemize}
\subsubsection{Montaje}
Para todos los materiales es el mismo procedimiento:
\begin{itemize}
    \item Una vez calibrada la balanza, se coloca un objeto de una clase en uno de los platos y otro objeto de la misma clase en el otro plato, ambos en las mismas condiciones; p.e., ambos globos desinflados, ambas bolsas vacías, o ambas botellas destapadas.
    \item Una vez que se haya llegado a la posición de equilibrio, se registran los resultados y se cambian las condiciones de uno de los objetos en la balanza para ver si cambia la posición de equilibrio.
    \item Se repite el experimento unas cuantas veces para corroborar los resultados obtenidos.
\end{itemize}

\subsection{Resultados}
Las imágenes muestran los resultados obtenidos:
\subsubsection{Globos}
\textbf{Primer experimento}

\begin{figure}[!h]
    \centering
    \includegraphics[scale=0.15]{/G11.jpeg}
    \caption{Primer experimento con globos, ambos desinflados.}
\end{figure}
\newpage
\begin{figure}[!h]
    \centering
    \includegraphics[scale=0.15]{/G12}
    \caption{Primer experimento con globos, con uno medianamente inflado.}
\end{figure}

\begin{figure}[!h]
    \centering
    \includegraphics[scale=0.15]{/G13.jpeg}
    \caption{Primer experimento con globos, con uno totalmente inflado.}
\end{figure}

\textbf{Segundo experimento}

\begin{figure}[!h]
    \centering
    \includegraphics[scale=0.15]{/G21.jpeg}
    \caption{Segundo experimento con globos, ambos desinflados.}
\end{figure}

\begin{figure}[!h]
    \centering
    \includegraphics[scale=0.15]{/G22.jpeg}
    \caption{Segundo experimento con globos, uno medianamente inflado.}
\end{figure}

\begin{figure}[!h]
    \centering
    \includegraphics[scale=0.15]{/G23.jpeg}
    \caption{Segundo experimento con globos, uno totalmente inflado.}
\end{figure}

\newpage
\subsubsection{Bolsas}

\textbf{Primer experimento}

\begin{figure}[!h]
    \centering
    \includegraphics[scale=0.15]{/B1.jpeg}
    \caption{Primer experimento con bolsas, ambas vacías.}
\end{figure}

\begin{figure}[!h]
    \centering
    \includegraphics[scale=0.15]{/B2.jpeg}
    \caption{Primer experimento con bolsas, una inflada.}
\end{figure}
\newpage
\textbf{Segundo experimento}

\begin{figure}[!h]
    \centering
    \includegraphics[scale=0.15]{/B21.jpeg}
    \caption{Segundo experimento con bolsas, ambas vacías.}
\end{figure}

\begin{figure}[!h]
    \centering
    \includegraphics[scale=0.15]{/B22.jpeg}
    \caption{Segundo experimento con bolsas, una inflada.}
\end{figure}
\newpage
\subsubsection{Botellas}

\textbf{Medianas}

\begin{figure}[!h]
    \centering
    \includegraphics[scale=0.15]{/BM1.jpeg}
    \caption{Primer experimento con botellas, ambas destapadas.}
\end{figure}

\begin{figure}[!h]
    \centering
    \includegraphics[scale=0.15]{/BM2.jpeg}
    \caption{Primer experimento con botellas, una tapada.}
\end{figure}
\newpage
\textbf{Segundo experimento}

\begin{figure}[!h]
    \centering
    \includegraphics[scale=0.15]{/BG1.jpeg}
    \caption{Segundo experimento con botellas, una tapada.}
\end{figure}

\begin{figure}[!h]
    \centering
    \includegraphics[scale=0.15]{/BG2.jpeg}
    \caption{Primer experimento con botellas, ambas destapadas.}
\end{figure}

\newpage
\subsection{Conclusiones}
En el caso de los globos, se nota una ligera diferencia en la lectura de la balanza; cuanto más se infla el globo, más se inclina hacia éste. En cambio, para las botellas y las bolsas, no sucede este cambio: todas las lecturas son iguales.

\subsection{Cuestionario}
\begin{itemize}
    \item \textbf{De acuerdo con la física involucrada en estos fenómenos, ¿qué explicación da usted al resultado de los globos?} 
    Como se explicó en la introducción, lo que mide la balanza es la fuerza normal del plato sobre el objeto, que es la fuerza de respuesta que contrarresta a la fuerza gravitacional: $F_{N}=F_{g}-F{f}$. Sin embargo, cuando se infla el globo, cambia esta ecuación, y tendremos:
    \begin{equation}
        F_{N'}=W_{G}+W{a}-V_{d}\rho_{f}\cdot g
    \end{equation}
    donde $W_{G}$ es el peso del globo y $W_{a}$ el peso del aire que contiene. Como las paredes del globo son elásticas, comprimen al fluido, aumentando la densidad del fluido. Por tanto, la fuerza normal de la balanza es mayor para el globo inflado que para el globo desinflado, y se inclina a ese lado.
    \item \textbf{¿Qué explicación da usted al resultado de las bolsas?}
    La explicación es similar a la de la pregunta anterior: a diferencia de las paredes del globo, que son elásticas, las paredes de la bolsa no son tan elásticas, por lo que la consiguiente compresión del fluido y aumento de la densidad son considerablemente menores, al grado de ser despreciables, por lo que la fuerza normal para las dos bolsas, la inflada y la desinflada, es prácticamente la misma; por tanto, no se inclina la balanza.
    \item \textbf{¿Qué explicación da usted al resultado de las botellas?}
    La misma que a la del caso de las bolsas: menor elasticidad, que lleva a menor compresión, que lleva a menor densidad, que lleva a un cambio imperceptible de fuerza normal.
    \item \textbf{Si súbitamente se endureciera el hule que forma el globo inflado, ¿esperaría usted
que ahora el resultado fuera como el de las botellas o las bolsas? ¿Por qué?}
    Puesto que el aire ya estaría comprimido, y la densidad aumentada, el resultado sería el mismo; sólo cambiaría si se endurecieran las paredes \emph{antes} de comprimir el aire.
    \item \textbf{En cierto programa televisivo infantil de divulgación de la ciencia se realizó el
siguiente experimento con el objeto de determinar si el aire tiene peso: se colocaron
dos globos inflados en los extremos de un balancín, se pinchó un globo y se observó
el resultado. ¿Usted qué opina?}
   Es algo como lo que se realizó en la práctica; sin embargo, el resultado de la práctica es debido a más factores que sólo el peso del globo, por lo que relacionar directamente lo que sucedió en el programa de televisión con la existencia del peso del aire no sería recomendable.
    \item \textbf{En la pregunta anterior, ¿Qué ocurriría si en lugar de usar globos se usaran bolsas de plástico y se sacara el aire de una de ellas?}
    Sucedería lo mismo que en la práctica: no de detectaría cambio alguno en la lectura de la balanza.
    \item \textbf{De acuerdo con su respuesta a la pregunta anterior, ¿qué podría concluir un observador acerca del peso del aire? Explique.}
    Que el peso (es decir, la fuerza de reacción a la fuerza de atracción gravitacional) del aire no sólo depende de su masa, sino del empuje; para recipientes más elásticos, mayor será el "peso" del aire.
    \item \textbf{¿Qué diría usted de la siguiente afirmación: “Se equilibran en una balanza una
    pelota de unicel y una bolita de plomo. Si ahora se introduce dicho sistema en una
    campana de vacío y se le extrae el aire se perderá el equilibrio”?. Argumente con ecuaciones su respuesta.}
    Tenemos inicialmente que 
   \begin{equation} \label{eq:4}
       F_{NU}=F_{NP}
   \end{equation}
         \begin{center}
      
        $\Rightarrow F_{gU}-F_{fU}=F_{gP}-F_{fP}$
       
       $\Rightarrow g(V_{U}(\rho_{U}-\rho_{aire}))=g(V_{P}(\rho_{P}-\rho_{aire}))$
    \end{center}
    Como $\rho_{P} > \rho_{U}$:
    \begin{center}
        $(\rho_{U}-\rho_{aire})<(\rho_{P}-\rho_{aire})$
        
        $\Rightarrow V_{U}>V_{P}$
    
        $\Rightarrow F_{fU}=V_{U}\rho_{aire}\cdot g>F_{fP}=V_{P}\rho_{aire}\cdot g$
    \end{center}
    Sumamos a ambos lados $F_{NU}$, y recordando de \ref{eq:2} que $mg=F_{N}+F_{f}$ y de \ref{eq:4} que $F_{NU}=F_{NP}$ obtenemos:
    \begin{center}
        $m_{U}\cdot g>m_{P}\cdot g$
        
        $m_{U}>m_{P}$
    \end{center}
    Es decir, cuando se use la campana de vacío y desaparezca el empuje, la balanza se inclinará hacia el lado de la bolita de unicel.
    
    \item \textbf{¿Qué le ha demostrado la solución de esta práctica?}
    Que para objetos ligeros, el empuje del aire, usualmente despreciable por su mínima contribución, se vuelve un factor que debe considerarse para mediciones precisas de peso.
    
    \item \textbf{¿Tiene usted razones o sugerencias para desechar o mejorar esta práctica en este laboratorio?} 
    No, puesto que no tuvimos problema con ella.
\end{itemize}
\bibliographystyle{ieeetr}
\bibliography{pda}
\end{document}

\end{document}